\section{Thảo luận}
\label{sec:discussion}

\subsection{Ba phát hiện chính}

Thứ nhất, ResNet-1D--TCN đem lại lợi ích vận hành rõ rệt mà điểm số tổng thể không phản ánh hết. Trên benchmark đã đo, thời gian suy luận giảm tương ứng mức tăng tốc 3,76 lần; chiến dịch huấn luyện và xác thực seed 123 của E3 hoàn tất trong 3 giờ 16 phút 40 giây, so với 33 giờ 35 phút 36 giây của E0. Những lợi ích này đi cùng số tham số cao hơn 4,37 lần và bộ nhớ GPU cực đại cao hơn 28,4\%. Đây là một đánh đổi có thể cân nhắc khi xây dựng pipeline, không phải sự thu nhỏ mô hình trên mọi phương diện.

Thứ hai, lợi ích dự báo phụ thuộc vào cấu hình và miền đánh giá. E1--E0 giữ encoder/đặc trưng nhưng thay cả mô hình chuỗi và lịch huấn luyện; E2--E1 thay gói đặc trưng/ngữ cảnh cùng cách huấn luyện encoder. Hai chênh lệch nhỏ trên Sleep-EDF chưa đạt ý nghĩa Wilcoxon sau Holm. Trên SHHS1, E1--E0 dương nhưng nhỏ, còn E2--E1 âm. Ngược lại, E3 vượt E0 và E6 trong hai đối chiếu SHHS chính đã khóa, đồng thời có chênh lệch gộp E3--E6 dương ở cả hai seed Sleep-EDF. Kết quả cho phép đánh giá các cấu hình đã triển khai, nhưng không tách một tác động kiến trúc thuần.

Thứ ba, lợi thế tổng thể không đồng nghĩa khắc phục được mọi lớp. E0 và E3 chỉ đạt recall N3 0,2610 và 0,2582 trên SHHS1, với hơn 72\% epoch N3 bị gán thành N2 ở cả hai mô hình. E6 còn có recall N3 thấp hơn, 0,2005. Sự tái diễn này là phát hiện trung tâm của phân tích chuyển cohort: pipeline E3 cải thiện điểm số và nhiều lớp khác, nhưng lỗi thiếu N3 vẫn tồn tại qua hai họ mô hình và không được pipeline z-score đã thử khắc phục.

\subsection{Đọc đúng các đối chiếu trong miền}

E3--E6 ở seed 42 là đối chiếu duy nhất đồng thời có CI bootstrap của chênh lệch gộp hoàn toàn dương và Wilcoxon có ý nghĩa sau Holm. Với E3--E2, CI gộp $[0,000305;0,014520]$ cũng dương, nhưng chỉ 37/78 người cải thiện và $p_{Holm}=0,8989$. Hai kết quả không mâu thuẫn: Macro-F1 từ tổng ma trận nhầm lẫn là hàm phi tuyến, còn Wilcoxon dùng các chênh lệch theo người. Chúng không thiết lập lợi ích đồng đều ở cấp đối tượng và không tự cho biết người có nhiều epoch hơn có tạo ra phần tăng gộp hay không.

Seed 123 giữ chênh lệch gộp E3--E6 dương với CI không chứa 0, nhưng không giữ ý nghĩa Wilcoxon sau Holm. Hướng hiệu ứng được lặp lại, trong khi độ lớn giảm từ 0,0213 xuống 0,0102. E4 cũng nhỉnh hơn E3 về Macro-F1 gộp ở seed 123. Vì vậy, E3 cao nhất ở chiến dịch seed 42, không phải đứng đầu mọi lần chạy; hai seed cố định chưa đủ mô tả đầy đủ biến thiên do khởi tạo.

\subsection{Lợi ích vận hành và phạm vi benchmark}

Việc giảm thời gian hoàn tất cross-validation từ hơn một ngày xuống vài giờ trên phần cứng đã đo rút ngắn vòng kiểm tra và phát triển mô hình. Tuy nhiên, wall-clock còn phụ thuộc số epoch thực chạy, optimizer, batch size, dừng sớm và việc E1 tái sử dụng encoder/cache E0. Không thể gán toàn bộ mức giảm cho phép thay BiLSTM bằng TCN. Tương tự, mức tăng tốc suy luận 3,76 lần là phép đo forward theo cửa sổ benchmark; nó không bao gồm tiền xử lý, đọc dữ liệu hoặc thời gian phải chờ ngữ cảnh tương lai. Lọc không lệch pha, BiLSTM hai chiều và TCN không nhân quả đặt hệ thống này trong bối cảnh phân tích ngoại tuyến.

\subsection{Tiền xử lý và đối chứng E6}

Trong đối chiếu hậu nghiệm SHHS1 đã quan sát, chênh lệch E3--E2 là $+0,04750$ Macro-F1 trung bình theo người, lớn hơn chênh lệch của E1--E0 và E2--E1. Phần mở rộng seed 123 cũng cho thấy E4 chỉ lọc dải cao hơn đầu vào thô E2, và cao hơn E3 cùng seed. Trong mẫu này, xử lý tín hiệu là một trục phát triển có ảnh hưởng đáng kể, không phải chi tiết triển khai trung tính. Một thiết kế nhân tố độc lập sẽ cần thiết để phân tách đóng góp của lọc dải, đổi thang và tương tác với huấn luyện.

E6 kiểm tra một lựa chọn cụ thể thường dễ áp dụng: z-score theo từng bản ghi. Trung bình và độ lệch chuẩn lấy từ toàn bộ recording đích không nhãn nên phép chuẩn hóa có tính transductive, dù trọng số vẫn chỉ học từ Sleep-EDF. Recall N3 toàn mẫu của E6 là 0,2005 so với 0,2582 của E3; tại vùng lân cận thay đổi nhãn, các giá trị là 0,0721 và 0,0733. Như vậy, pipeline z-score đã thử không khắc phục thiếu N3. Do E6 được huấn luyện riêng với đầu vào chuẩn hóa tương ứng, đây không phải can thiệp chỉ thay biên độ ở cùng một mô hình; kết quả không loại trừ vai trò của biên độ, montage hoặc hiệu quả của các phương pháp thích nghi không nhãn khác.

\subsection{Ý nghĩa của lỗi N3 và phân tích phản thực}

E3 có precision N3 0,9440 nhưng recall 0,2582: khi phát lớp N3 mô hình thường đúng, song bỏ sót phần lớn N3 thật. E0 có mẫu thiếu N3 gần tương tự. Cả hai có recall N3 thấp hơn gần thay đổi nhãn tham chiếu so với vùng ổn định, nhưng ngay trong vùng ổn định SHHS1 recall cũng chỉ là 0,4008/0,3900. Vì vậy, mô tả lỗi chỉ quanh ranh giới không bao quát toàn bộ thiếu hụt. Thành phần lớp giữa các vùng khác nhau; chênh lệch metric không xác định rằng chuyển pha sinh lý gây ra lỗi. Montage, thiết bị, quần thể và thực hành chấm cũng thay đổi đồng thời giữa hai dataset, nên chưa thể quy lỗi cho một yếu tố riêng.

Oracle sửa riêng N3$\rightarrow$N2 trong ma trận E3 nâng Macro-F1 SHHS1 từ 0,6099 lên 0,7444, tương ứng 74,5\% chênh lệch điểm quan sát với Sleep-EDF. Đây là cách định lượng quy mô kênh lỗi, không phải hiệu năng một phương pháp sẵn có hoặc phần mất mát nhân quả do dịch chuyển miền. Sleep-EDF dùng dự đoán ngoài fold, còn SHHS dùng ensemble mười mô hình; phân bố lớp và thành phần cohort cũng khác. N3$\rightarrow$N2, tiếp theo là N2$\rightarrow$REM trong E3, là các mục tiêu hợp lý cho một nghiên cứu khắc phục có kiểm chứng, nhưng oracle chưa chứng minh chiến lược thu nhãn nào tối ưu.

\subsection{Vai trò của ngữ cảnh và không gian đặc trưng}

Gate 8 chưa phát hiện lợi ích tăng thêm có ý nghĩa của các nhóm P/N được mã hóa riêng trong các điều kiện huấn luyện lại TCN. Kết quả không thiết lập tương đương hay phủ định giá trị của ngữ cảnh thời gian nói chung. Silhouette E2 thấp hơn E1 ở cả mười fold cũng chỉ mô tả cấu trúc cụm theo khoảng cách Euclid sau pipeline phân tích; không đủ xếp hạng khả năng dự đoán chuỗi của hai biểu diễn. Các phân tích này giúp giới hạn những cách giải thích về biểu diễn và ngữ cảnh, không thay thế bằng chứng hiệu năng bắt cặp.

\subsection{Giới hạn và nghiên cứu tiếp theo}

Nghiên cứu chỉ khảo sát một hướng chuyển Sleep-EDF sang mẫu 180 người SHHS1, hai seed cố định và các cấu hình cụ thể. Một số đối chiếu thành phần, phần mở rộng E4 và vùng nhãn là hậu nghiệm trên cohort đã xem kết quả. Cùng subject splits bảo vệ tính bắt cặp nhưng không biến các training recipe khác nhau thành đối chứng kiến trúc thuần. Các số liệu này hỗ trợ nghiên cứu thực nghiệm về pipeline và lỗi chuyển cohort, chưa xác nhận hiệu quả lâm sàng.

Bước kiểm chứng phù hợp có thể là lặp lại mẫu lỗi trên một cohort chưa được xem hoặc trên một pipeline đối chứng mạnh dưới cùng giao thức. Nếu mục tiêu chuyển sang khắc phục N3, có thể đánh giá calibration hoặc fine-tuning trên tập thích nghi/validation riêng, báo cáo đồng thời recall, precision và false-positive N3, ảnh hưởng lên N2, cùng kết quả tại vùng lân cận và vùng ổn định. Không dùng 180 test subjects để chọn threshold hoặc cập nhật mô hình. Nghiên cứu hiện tại chưa so sánh chi phí thu nhãn hay hiệu quả của các hướng này, nên không xếp chúng thành một thứ tự tối ưu về công sức.
